\section{Analysis}\label{sec:analysis}

The four-way ablation in \S\ref{sec:experiments} established that gating produces no measurable benefit over LoRA-only fine-tuning. This section investigates why. We first characterize the gradient dynamics between the two training objectives (\S\ref{subsec:gradient-conflict}), then show that the learned gates exhibit behavior opposite to the design intent (\S\ref{subsec:gate-conditioning}). A targeted causal experiment rules out gradient direction conflict as the explanation (\S\ref{subsec:pcgrad}), and a distribution shift analysis identifies the actual failure mode (\S\ref{subsec:dist-shift}).

\subsection{Gradient Conflict Observation}\label{subsec:gradient-conflict}

To understand how the two loss terms interact during training, we measure the cosine similarity between $\nabla_\theta \mathcal{L}_\text{NTP}$ and $\nabla_\theta \mathcal{L}_\text{GAR}$ at each transformer layer across three training epochs. Table~\ref{tab:grad-conflict} summarizes the results.

\begin{table}[t]
\centering
\small
\caption{Gradient interaction between $\mathcal{L}_\text{NTP}$ and $\mathcal{L}_\text{GAR}$ across training epochs. \textit{cos\_sim}: mean cosine similarity across layers. \textit{neg\%}: percentage of layers with negative cosine similarity. \textit{mag\_ratio}: ratio of NTP to GAR gradient magnitudes.}
\label{tab:grad-conflict}
\begin{tabular}{lccc}
\toprule
\textbf{Epoch} & \textbf{cos\_sim} & \textbf{neg\%} & \textbf{mag\_ratio} \\
\midrule
0 & $-$0.089 & 78.6 & 7.5$\times$ \\
1 & $-$0.054 & 71.4 & 13.7$\times$ \\
2 & $-$0.042 & 67.9 & 14.7$\times$ \\
\bottomrule
\end{tabular}
\end{table}

Two patterns are apparent. First, the NTP and GAR gradients point in opposing directions across 68--79\% of layers throughout training. The mean cosine similarity is negative at every epoch, though it attenuates slightly from $-$0.089 to $-$0.042 as training progresses. Second, the NTP gradient magnitude dominates the GAR gradient by a factor that grows from 7.5$\times$ at epoch~0 to 14.7$\times$ at epoch~2. Direction conflict decreases while magnitude imbalance worsens. We observe these patterns as descriptive statistics of the training dynamics. This does not establish that gradient conflict causes $\mathcal{L}_\text{GAR}$ to fail; \S\ref{subsec:pcgrad} tests this hypothesis directly.

\subsection{Gate Turn Conditioning}\label{subsec:gate-conditioning}

The design intent of ThreadGate is to modulate the residual stream more aggressively at late conversation turns, where GAR naturally decays (\S\ref{sec:background}). We test whether the learned gates achieve this by comparing gate values at turn~1 (early) and turn~21 (late) across all 28 layers of the trained model.

The gates do the opposite. In 27 of 28 layers, the gate value at turn~1 is lower than at turn~21 (Wilcoxon signed-rank test: $p = 3.73 \times 10^{-8}$, effect size $r = 0.975$). At layer~23, the gate is 0.769 at turn~1 and 0.979 at turn~21; at layer~14, it is 0.747 at turn~1 and 0.928 at turn~21. The pattern is consistent: gates suppress the residual stream more at early turns (where GAR is already high) and become nearly passive at late turns (where intervention is most needed). This inversion has a straightforward explanation. NTP loss pressure is highest at late turns, where the model must predict tokens conditioned on a long context. The NTP gradient pushes gates toward 1.0 to preserve the full sublayer computation, and this pressure is strongest precisely where GAR-based attenuation would be most useful.

The practical consequence is that the gates provide almost no intervention at inference time. Across all layers, mean gate values during inference range from 0.93 to 0.98, corresponding to residual stream attenuation of only 2--7\%. This attenuation is far below what would be required for meaningful modulation, and it explains the convergence with LoRA-only in Table~\ref{tab:four-way}. Figure~\ref{fig:gate-shift}(a) visualizes the per-layer gate difference between early and late turns.

\subsection{PCGrad Causal Exclusion}\label{subsec:pcgrad}

The gradient conflict observed in \S\ref{subsec:gradient-conflict} raises a natural hypothesis: perhaps the opposing gradient directions between $\mathcal{L}_\text{NTP}$ and $\mathcal{L}_\text{GAR}$ prevent the GAR objective from converging. If direction conflict is the cause of $\mathcal{L}_\text{GAR}$ collapse, then eliminating direction conflict should prevent the collapse. We test this prediction using PCGrad \citep{yu2020pcgrad}.

PCGrad projects the gradient of one task onto the normal plane of the other whenever two task gradients conflict in direction. We apply PCGrad to the gate parameters only, projecting $\nabla_\theta \mathcal{L}_\text{NTP}$ onto the orthogonal complement of $\nabla_\theta \mathcal{L}_\text{GAR}$ whenever their cosine similarity is negative. This eliminates the directional component of the conflict while leaving magnitude differences intact. PCGrad does not address magnitude imbalance (which GradNorm targets; \citealp{chen2018gradnorm}) or Nash equilibrium violations (which Nash-MTL targets; \citealp{navon2022nash}); our claim is scoped to direction conflict alone.

The result is unambiguous. Under PCGrad, the cosine similarity between NTP and GAR gradients on gate parameters turns positive by gradient step~200 ($+$0.086), confirming that PCGrad successfully removes direction conflict. Yet $\mathcal{L}_\text{GAR}$ follows a trajectory nearly identical to the baseline (v7) and collapses in the same manner. Figure~\ref{fig:training-pcgrad}(b) overlays the $\mathcal{L}_\text{GAR}$ training curves for both conditions.

The diagnostic exclusion follows by \emph{modus tollens}:
\begin{enumerate}
    \item If direction conflict causes $\mathcal{L}_\text{GAR}$ collapse, then removing direction conflict should prevent the collapse.
    \item PCGrad removes direction conflict (cosine similarity turns positive at step~200).
    \item $\mathcal{L}_\text{GAR}$ collapses regardless.
    \item Therefore, direction conflict is not the cause of $\mathcal{L}_\text{GAR}$ collapse.
\end{enumerate}

This result recharacterizes the gradient observations in \S\ref{subsec:gradient-conflict}: direction conflict between $\mathcal{L}_\text{NTP}$ and $\mathcal{L}_\text{GAR}$ is a symptom of the multi-objective optimization dynamics, not the mechanism responsible for GAR failure. The magnitude imbalance (7.5--14.7$\times$) and potential Nash equilibrium violations remain as candidate explanations that future work can test with GradNorm \citep{chen2018gradnorm} and Nash-MTL \citep{navon2022nash}.

\subsection{Distribution Shift Analysis}\label{subsec:dist-shift}

The preceding subsections explain why the gating mechanism is ineffective during training. This subsection addresses why the learned GAR improvements fail to transfer from training to inference.

Recall from \S\ref{sec:experiments} that teacher-forced GAR-ours reaches 0.49 for all fine-tuned variants, while autoregressive GAR-ours drops to 0.016 for ThreadGate, a 30$\times$ gap. We decompose this gap by comparing same-distribution and cross-distribution evaluation. Under same-distribution evaluation (teacher-forced UltraChat text), GAR-ours is approximately 0.49, confirming that the model has learned the intended attention pattern on the training distribution. Under cross-distribution evaluation (model-generated text), GAR-ours drops to 0.016. The training itself is sound; the failure is one of generalization.

GAR-Dongre provides a revealing control. Across both evaluation conditions, GAR-Dongre remains uniformly low (0.01--0.02) with no distribution gap. The distribution shift selectively destroys the concentrated attention pattern captured by GAR-ours (last-token, summary-position attention) while leaving the diffuse attention pattern captured by GAR-Dongre (averaged across all positions) unaffected. This selectivity implies that the distribution shift disrupts a specific learned routing behavior rather than attention computation in general.

The mechanism connecting training regime to distribution shift is exposure bias \citep{bengio2015scheduled}. During training, the model conditions each token prediction on ground-truth prefixes from UltraChat. At inference, it conditions on its own prior outputs. The attention patterns optimized for the former distribution do not transfer to the latter, because the statistical properties of model-generated text diverge from those of human-written reference text. The 2--7\% gate attenuation documented in \S\ref{subsec:gate-conditioning} is far too weak to compensate for this distributional mismatch.

\citet{he2019exposure} showed that language models exhibit token-level self-recovery from exposure bias: accumulated errors in autoregressive generation do not degrade perplexity as severely as worst-case analysis predicts. Our results are consistent with this finding at the token level, as PPL remains normal across all conditions (Table~\ref{tab:four-way}). The critical distinction is between primary and auxiliary objectives. The NTP objective is robust to exposure bias because next-token prediction is the same task at both training and inference time. The GAR auxiliary objective is fragile because it optimizes an attention pattern under teacher-forced context that has no counterpart in autoregressive generation. Primary-task self-recovery does not extend to auxiliary objectives that depend on distributional properties of the input prefix. Figure~\ref{fig:gate-shift}(b) summarizes the four-model GAR comparison across evaluation conditions.